\chapter{Conclusion and Future Work}

\section{Conclusion}

Cloud-native systems have become the foundation of modern software services, yet their security remains fundamentally challenged by software vulnerabilities in cloud infrastructure, external attacks under dynamic system behaviors, and human-driven insider adversarial activities. In cloud-native environments, external attacks may originate from exploitable software vulnerabilities and further evolve across multiple layers of the system, while insider threats arise from privileged or trusted entities operating within organizational boundaries. This dissertation improves the current understanding of these security problems through a progressive and threat-driven perspective, moving from infrastructure-level defect analysis to runtime attack detection and, ultimately, to insider threat analysis.

The main contributions of this dissertation are summarized as follows:

\begin{itemize}

\item \textbf{Defect Characterization in Cloud Infrastructure.} This dissertation presents a systematic empirical study of defects in container runtime systems. By analyzing real-world bugs, it identifies recurring defect symptom patterns and root causes, and reveals fundamental limitations of existing testing techniques in detecting these defects. These findings provide actionable insights for improving testing practices and reducing exploitable vulnerabilities in critical cloud components.

\item \textbf{Robust Detection of External Attacks under Dynamic Environments.} This dissertation investigates the effectiveness of log-based attack detection methods in non-stationary cloud environments. It introduces CAShift, a benchmark that models realistic system shifts caused by application updates, software evolution, and infrastructure changes. The results demonstrate significant performance degradation of existing methods under such conditions, and further provide a systematic adoption of continual learning strategies for shift adaptation, offering insights into their effectiveness in improving detection robustness under evolving system behaviors.

\item \textbf{Data-Driven Insider Threat Detection.} This dissertation proposes Chimera, an LLM-based multi-agent framework for simulating behaviorally diverse organizational activities. It models individuals as interactive agents to generate complex, human-driven behaviors. Based on this framework, it constructs ChimeraLog, a large-scale multi-modal dataset capturing both adversarial insider behaviors and high-fidelity benign activities. The generated data is validated through human expert evaluation and statistical analysis, demonstrating its realism and diversity, and highlighting the importance of data realism in enabling effective insider threat analysis.

\end{itemize}

\section{Future Work}

Building upon the findings of this dissertation, it identifies promising directions that can be explored to further advance the trustworthiness and security of cloud-native systems.